\documentclass[11pt]{article}
\usepackage[margin=1in]{geometry}
\usepackage{amsmath,amssymb,amsthm,mathtools}
\usepackage[T1]{fontenc}
\usepackage{lmodern,microtype,hyperref}
\hypersetup{colorlinks=true,linkcolor=blue,citecolor=blue,urlcolor=blue,
 pdftitle={A logarithmic refinement of CKN partial regularity},
 pdfauthor={Yongxi Lin}}
\setlength{\emergencystretch}{2em}
\newtheorem{theorem}{Theorem}
\newtheorem{lemma}[theorem]{Lemma}
\newtheorem{proposition}[theorem]{Proposition}
\newcommand{\R}{\mathbb R}
\newcommand{\esssup}{\operatorname*{ess\,sup}}
\newcommand{\eps}{\varepsilon_*}
\title{A logarithmic refinement of CKN partial regularity}
\author{Yongxi Lin}
\date{}
\begin{document}
\maketitle

\begin{abstract}
Let \((u,p)\) be a suitable weak solution of the three-dimensional
Navier--Stokes equations.
We prove that its singular set has zero parabolic Hausdorff
measure for the gauge
$h(r)=r[\log(1/r)]^2.$
The proof is generated by AI under my guidance.
\end{abstract}

\section{Introduction}

Consider the Navier--Stokes equations with viscosity normalized to one:
$$
\partial_tu-\Delta u+(u\cdot\nabla)u+\nabla p=0,
\qquad \nabla\cdot u=0.
$$
Caffarelli, Kohn, and Nirenberg \cite{CKN} proved that the interior
singular set of a suitable weak solution has zero one-dimensional
parabolic Hausdorff measure. Logarithmic refinements of this conclusion
have also been studied; see, for example, Lei and Ren \cite{LR}.
This article explains an argument giving a squared logarithmic
factor. This factor strengthens the measure conclusion while leaving
the critical power of \(r\) equal to one.

Let \(\mathcal D\subset\R^3\times\R\) be open. Denote by \(S\) the
interior singular set: the points where \(u\) is not essentially bounded
in any open spacetime neighborhood.

Write
$$
C_r(x,t)=B_r(x)\times(t-r^2,t+r^2).
$$
For an increasing gauge \(h\), define the parabolic Hausdorff measure by
$$
\mathcal P^h(A)=\lim_{\delta\downarrow0}
\inf\left\{\sum_i h(r_i):
 A\subset\bigcup_i C_{r_i}(z_i),\quad 0<r_i<\delta\right\}.
$$
This is equivalent, up to fixed constants for the gauge used below,
to Hausdorff measure in the parabolic metric
\(d((x,t),(y,s))=\max\{|x-y|,|t-s|^{1/2}\}\).
Changing between radii and diameters therefore does not affect nullity.

\begin{theorem}\label{main}
Let \((u,p)\) be a suitable weak solution on \(\mathcal D\). Then
$$
\mathcal P^h(S)=0,
\qquad h(r)=r[\log(1/r)]^2.
$$
The gauge is specified for sufficiently small \(r\) and extended
increasingly to larger radii.
\end{theorem}

\paragraph{Main idea and physical interpretation.}
The local energy inequality compares the change in kinetic energy and
the viscous dissipation with energy transported through the boundary
of a cylinder. The pressure contributes to this transport. On a spatial
scale \(r\), the diffusion time is \(r^2\), which explains the parabolic
cylinders used here. On each compact localization, the finite measure
\(|\nabla u|^2\,dx\,dt\) records the viscous dissipation.

At a singular point, the scaled cubic velocity and pressure integral
cannot become small. A local estimate says that this activity contracts
at a smaller scale unless a certain mixed gradient integral pays for
it. We use spatial exponent \(12/7\) because Sobolev embedding then
gives \(L^4\), and interpolation with \(L^2\) matches the cubic velocity
term. The resulting cost contains the square root of the activity.
The use of this mixed gradient exponent is consistent with earlier
regularity criteria; see \cite{GKT}.

A one-scale covering argument does not compare the costs at different
scales. Our trace estimate does: it sums the mixed dissipation against
a measure supported on a hypothetically large singular set. The squared
logarithmic factor corresponds to the weights \(1/j\). Their
sum diverges, while the trace estimate gives a finite integral of the
same weighted cost. The logarithmic factor arises from this summation over
scales, rather than from an additional physical assumption.

The proof has three steps.
\begin{enumerate}
\item Obtain a contraction estimate and show that persistent activity
has infinite weighted dissipation cost.
\item Prove a trace estimate that bounds the total cost using the
finite viscous dissipation.
\item Apply Frostman's lemma and compare cylinders with dyadic boxes
 to obtain a contradiction.
\end{enumerate}

\section{The cost of persistent activity}

We use the standard notion of a suitable weak solution: the equations
hold in distributions, locally
$$
u\in L^\infty_tL^2_x\cap L^2_tH^1_x,
\qquad p\in L^{3/2}_{x,t},
$$
and the local energy inequality holds.

All estimates are local. After a fixed translation and rescaling,
work on \(B_4\times(-4,4)\Subset\mathcal D\), with centers in a compact
set \(L\subset B_1\times(-1,1)\). Choose all radii sufficiently small
that the cylinders used below lie in \(U=B_2\times(-2,2)\).
A countable family of such localizations covers \(\mathcal D\).
Fixed rescaling preserves the nullity conclusion because
\(h(cr)\asymp h(r)\) for each fixed \(c>0\).

For \(z=(x,t)\), define
$$
\begin{aligned}
G(r,z)&=r^{-2}\int_{C_r(z)}(|u|^3+|p|^{3/2})\,dy\,ds,\\
A(r,z)&=r^{-1}\esssup_{|s-t|<r^2}
                  \int_{B_r(x)}|u(y,s)|^2\,dy,\\
E(r,z)&=r^{-1}\int_{C_r(z)}|\nabla u|^2\,dy\,ds,\\
e(r,z)&=r^{-3/2}\int_{t-r^2}^{t+r^2}
       \left(\int_{B_r(x)}|\nabla u(y,s)|^{12/7}\,dy\right)^{7/6}ds.
\end{aligned}
$$
These quantities are invariant under Navier--Stokes scaling. Spatial
H\"older gives \(e(r,z)\leq CE(r,z)\).
A standard cutoff in the local energy inequality gives
$$
A(r/2,z)+E(r/2,z)\leq C\bigl(G(r,z)^{2/3}+G(r,z)\bigr).
$$
Here the quadratic cutoff term gives \(G^{2/3}\), and the cubic and
pressure flux terms give \(G\); for the latter use H\"older with the
pressure exponent \(3/2\). See also the local energy estimates in
\cite{GP}.

The standard velocity--pressure epsilon-regularity criterion implies
that, for a fixed \(\eps>0\),
$$
G(r,z)\geq\eps
\quad\text{at every }z\in S\cap L\text{ and every small }r.
$$
To apply the backward-cylinder formulation to our symmetric cylinders,
use the backward cylinder of radius \(r\) with top time \(t+r^2/8\).
It lies in \(C_r(z)\), and its smaller regularity cylinder contains an
open neighborhood of \(z\). Thus smallness of \(G(r,z)\) implies
regularity at \(z\). The precise criterion is recorded in
\cite[Theorem 1.1]{GP}.

\begin{lemma}[Scale recurrence]\label{recurrence}
There are fixed constants \(K=512\), \(\theta=2^{-M}<1/(2K)\), and
\(C_*>0\) such that
$$
G(\theta r,z)\leq\tfrac14G(r,z)
   +C_*\sqrt{G(r,z)}\,e(r/K,z)
\qquad\text{whenever }G(r,z)\geq\eps.
$$
\end{lemma}

\begin{proof}
We first show that for \(0<s\leq\rho/2\), suppressing the center,
$$
G(s)\leq C\frac{s}{\rho}G(\rho)
       +C\left(\frac{\rho}{s}\right)^2 A(\rho)^{1/2}e(\rho).
$$
At each time, let \(b\) be the spatial mean of \(u\) on \(B_\rho\).
Sobolev--Poincar\'e and interpolation give
$$
\|u-b\|_4\leq C\|\nabla u\|_{12/7},
\qquad \|u-b\|_3^3\leq\|u-b\|_2\|u-b\|_4^2.
$$
All norms here are over \(B_\rho\). There is no radius factor in the
first inequality because \(1/(12/7)-1/4=1/3\).
Taking the \(L^2\) norm uniformly in time and integrating the other
factor gives
$$
\int_{C_\rho}|u-b|^3
\leq C\left(\esssup_s\int_{B_\rho}|u|^2\right)^{1/2}
 \int_{t-\rho^2}^{t+\rho^2}
       \left(\int_{B_\rho}|\nabla u|^{12/7}\right)^{7/6}ds
=C\rho^2\sqrt{A(\rho)}\,e(\rho).
$$
Jensen's inequality bounds the mean contribution on \(C_s\), after
division by \(s^2\), by \(C(s/\rho)G(\rho)\).

For the pressure, decompose at each time
$$
p_1=R_iR_j[\mathbf1_{B_\rho}(u_i-b_i)(u_j-b_j)],
\qquad p_2=p-p_1,
$$
where \(R_i\) are the Riesz transforms and repeated indices are summed.
The pressure equation and incompressibility imply that \(p_2\) is
harmonic on \(B_\rho\). The standard estimates give
$$
\int_{\R^3}|p_1|^{3/2}\leq C\int_{B_\rho}|u-b|^3,
\qquad
\int_{B_s}|p_2|^{3/2}\leq C(s/\rho)^3\int_{B_\rho}|p_2|^{3/2}.
$$
Integrating over the smaller time interval and using the preceding
velocity bound proves the claimed estimate for \(G(s)\).

Now set \(\rho=r/K\) and \(s=\theta r\). Since
$$
G(r/K)\leq K^2G(r),
\qquad A(r/K)\leq(K/2)A(r/2),
$$
the local energy bound gives
$$
G(\theta r)\leq CK^3\theta G(r)
 +C_{K,\theta}\bigl(G(r)^{1/3}+G(r)^{1/2}\bigr)e(r/K).
$$
Choose a dyadic \(\theta\) so that \(CK^3\theta\leq1/4\).
For \(G(r)\geq\eps\), the factor \(G(r)^{1/3}\) is bounded by a
constant times \(G(r)^{1/2}\). This proves the recurrence.
\end{proof}

\begin{lemma}[Persistent activity]\label{divergence}
Put \(r_j=R\theta^j\), where \(R\) is small, and let
\(G_j(z)=G(r_j,z)\), \(e_j(z)=e(r_j/K,z)\).
If \(G_j(z)\geq\eps\) for every \(j\), then
$$
\sum_{j\geq2}\frac{e_j(z)}{j\sqrt{G_j(z)}}=\infty.
$$
\end{lemma}

\begin{proof}
Lemma \ref{recurrence} and \(\log(1+s)\leq s\) give
$$
\log(G_{j+1}/\eps)-\log(G_j/\eps)
\leq-\log4+4C_*\frac{e_j}{\sqrt{G_j}}.
$$
Multiply by \(1/j\) and sum. Since \(G_j\geq\eps\),
summation by parts gives
$$
\begin{aligned}
\sum_{j=2}^N\frac{\log(G_{j+1}/\eps)-\log(G_j/\eps)}{j}
&=\frac{\log(G_{N+1}/\eps)}{N}
 -\frac12\log(G_2/\eps)\\
&\quad+\sum_{j=3}^N\frac{\log(G_j/\eps)}{j(j-1)}\\
&\geq-\frac12\log(G_2/\eps).
\end{aligned}
$$
Therefore
$$
4C_*\sum_{j=2}^N\frac{e_j}{j\sqrt{G_j}}
\geq(\log4)\sum_{j=2}^N\frac1j-\frac12\log(G_2/\eps)
\longrightarrow\infty.
$$
\end{proof}

\section{A trace estimate for the dissipation}

We now prove the estimate that sums the dissipation cost over all
scales. Define the finite measure
$$
\nu(A)=\int_{A\cap U}|\nabla u|^2\,dx\,dt.
$$
In the following formulas, \(|\nabla u|\) is extended by zero outside
\(U\); no extension of \(u\) is needed.
Fix a nested parabolic dyadic grid. A level-\(n\) box \(Q=X\times I\)
has spatial side \(\ell_n=2^{-n}\) and temporal length \(\ell_n^2\).
Use half-open boxes so that each level partitions spacetime. Define
$$
E_Q=\frac{\nu(Q)}{\ell_n},
\qquad
e_Q=\ell_n^{-3/2}\int_I
       \left(\int_X|\nabla u|^{12/7}\,dx\right)^{7/6}dt.
$$
These are the box versions of \(E\) and \(e\). A quotient with
\(\nu(Q)=0\) is set to zero; in that case \(e_Q=0\).

\begin{lemma}[Dissipation trace]\label{trace}
Let \(\mu\) be a finite nonnegative measure. Suppose that \(F(n)>0\)
is nondecreasing for \(n\geq n_0\) and
$$
\mu(Q)\leq C_\mu\ell_nF(n)
\quad\text{for every level-}n\text{ box},\quad n\geq n_0.
$$
Then
$$
\sum_{n\geq n_0}\sum_{\operatorname{level}(Q)=n}
 \mu(Q)\frac{e_Q}{\sqrt{E_QF(n)}}
\leq C\sqrt{\mu(\R^4)\nu(\R^4)}.
$$
The constant depends only on \(C_\mu\) and the spatial dimension.
\end{lemma}

\begin{proof}
The proof has two steps: a local Carleson estimate, followed by a
summation over the ratio of the two masses.

\paragraph{Step 1: an estimate inside each root.}
Put
$$
\beta_Q=\sqrt{\frac{\mu(Q)\ell_n}{F(n)}}\,e_Q.
$$
At a fixed time \(t\), let \(I_n(t)\) be its level-\(n\) grid interval.
For a level-\(n\) spatial cube \(X\), define
$$
\alpha_X(t)=\frac{\ell_n^{5/2}}{\sqrt{F(n)}}
                 \sqrt{\mu(X\times I_n(t))}.
$$
For a spatial ancestor \(X_0\) at level \(k\geq n_0\),
Cauchy--Schwarz over the \((\ell_k/\ell_n)^3\) spatial descendants
and nesting of the time intervals give
$$
\begin{aligned}
\sum_{X\subseteq X_0}\alpha_X(t)
&\leq\ell_k^{3/2}\sqrt{\mu(X_0\times I_k(t))}
          \sum_{n\geq k}\frac{\ell_n}{\sqrt{F(n)}}\\
&\leq2\frac{\ell_k^{5/2}}{\sqrt{F(k)}}
            \sqrt{\mu(X_0\times I_k(t))}
\leq C|X_0|.
\end{aligned}
$$
Thus the coefficients form a spatial Carleson sequence, uniformly in
\(t\). Apply the standard dyadic Carleson embedding at exponent
\(7/6\) to \(|\nabla u(\cdot,t)|^{12/7}\), restricted to \(X_0\).
Writing the spatial averages explicitly, it gives
$$
\sum_{X\subseteq X_0}\alpha_X(t)
  \left(\frac1{|X|}\int_X|\nabla u(x,t)|^{12/7}\,dx\right)^{7/6}
\leq C\int_{X_0}|\nabla u(x,t)|^2\,dx.
$$
Since
\(e_Q=\ell_n^2\int_I
\left(\frac1{|X|}\int_X|\nabla u(x,t)|^{12/7}\,dx\right)^{7/6}\,dt\),
integrating over the time interval of a root \(Q_0=X_0\times I_0\)
and using Tonelli proves
$$
\sum_{Q\subseteq Q_0}\beta_Q\leq C\nu(Q_0).
$$
The sum includes \(Q_0\) and its descendants. The essential point is
that this estimate holds inside every root.

\paragraph{Step 2: summation by the mass ratio.}
Fix \(\lambda>0\) and consider boxes with \(\nu(Q)>0\) and
\(\mu(Q)>\lambda\nu(Q)\). Their maximal ancestors satisfying the
same condition, within the level-\(n_0\) roots, are disjoint.
Applying Step 1 on these ancestors gives
$$
\sum_{\substack{\nu(Q)>0\\\mu(Q)>\lambda\nu(Q)}}\beta_Q
\leq C\min\left\{\nu(\R^4),\frac{\mu(\R^4)}{\lambda}\right\}.
$$
Layer cake therefore yields
$$
\begin{aligned}
\sum_{\nu(Q)>0}\beta_Q\sqrt{\frac{\mu(Q)}{\nu(Q)}}
&=\frac12\int_0^\infty\lambda^{-1/2}
 \sum_{\substack{\nu(Q)>0\\\mu(Q)>\lambda\nu(Q)}}\beta_Q\,d\lambda\\
&\leq2C\sqrt{\mu(\R^4)\nu(\R^4)}.
\end{aligned}
$$
The assertion is immediate if either total mass is zero. Finally,
$$
\beta_Q\sqrt{\frac{\mu(Q)}{\nu(Q)}}
=\mu(Q)\frac{e_Q}{\sqrt{E_QF(n)}},
$$
which proves the lemma.
\end{proof}

\section{Proof of Theorem \ref{main}}

Fix the compact localization \(L\) from Section 2. Choose a small dyadic
\(R\) and use \(r_j=R\theta^j\) as above. We argue by contradiction,
assuming \(\mathcal P^h(S\cap L)>0\).

\paragraph{Step 1: a measure on the persistent set.}
For each fixed \(r\), the function \(G(r,z)\) is continuous in \(z\),
as an integral of a locally integrable density over translated
cylinders. Thus
$$
T=\{z\in L:G(r_j,z)\geq\eps\text{ for every }j\geq0\}
$$
is compact and contains \(S\cap L\). Frostman's lemma for a doubling
gauge in the parabolic metric supplies a probability measure supported
on \(T\) such that
$$
\mu(C_r(z))\leq Ch(r)
\quad\text{for all centers and all sufficiently small }r.
$$
See \cite{BP} for the gauge formulation of Frostman's lemma.

\paragraph{Step 2: comparison of cylinders and boxes.}
Use the standard one-third trick: three adjacent nested grids in each
spatial coordinate and three in time give \(81\) parabolic product
grids. For a dyadic \(\rho\), every \(C_\rho(z)\) is contained in a
box with spatial side \(16\rho\) and temporal length \((16\rho)^2\)
in one of these grids. This is the usual adjacent-grid containment
property; see \cite{HK} for the general construction.
Put \(\rho=r_j/K\) and choose such a box \(Q\).
The box has level \(n_j=Mj+c\), where \(c\) is a fixed integer.
Moreover \(Q\subset C_{r_j/2}(z)\), since
$$
16\sqrt3\,\rho<r_j/2,
\qquad 256\rho^2<(r_j/2)^2.
$$
Because \(Q\) contains the open cylinder centered at \(z\), it is
that grid's unique level-\(n_j\) box containing \(z\).
Monotonicity of the unnormalized integrals and the local energy bound
now give, for \(z\in T\),
$$
e_j(z)\leq16^{3/2}e_Q,
\qquad
E_Q\leq16E(r_j/2,z)
\leq C\bigl(G_j(z)^{2/3}+G_j(z)\bigr)\leq C_*G_j(z).
$$
Consequently
$$
\frac{e_j(z)}{\sqrt{G_j(z)}}
\leq C\frac{e_Q}{\sqrt{E_Q}},
$$
with the zero convention already specified.

\paragraph{Step 3: finite trace versus infinite cost.}
Each fine box lies in a cylinder of comparable radius. The Frostman
bound therefore implies, in every one of the grids,
$$
\mu(Q)\leq C\ell_nF(n),\qquad F(n)=n^2,
$$
because \(h(C\ell_n)\asymp\ell_n n^2\).
Apply Lemma \ref{trace} and write \(Q_n(z)\) for the box containing
\(z\). It follows that
$$
\int_T\sum_{n\geq n_0}
 \frac{e_{Q_n(z)}}{n\sqrt{E_{Q_n(z)}}}\,d\mu(z)<\infty.
$$
Since \(n_j=Mj+c\), its weight is comparable to \(1/j\).
The comparison in Step 2 and summation over all \(81\) grids thus give
$$
\int_T\sum_{j\geq j_0}
 \frac{e_j(z)}{j\sqrt{G_j(z)}}\,d\mu(z)<\infty.
$$
Using the sum over the grids avoids having to choose a grid measurably
for each \((z,j)\). But Lemma \ref{divergence} makes the integrand
infinite at every point of \(T\), and \(\mu(T)=1\). This contradiction
proves \(\mathcal P^h(S\cap L)=0\). A countable localization completes
the proof. \hfill\(\square\)

\section{A short extension}

The preceding proof uses only the harmonic weights \(1/j\).
Inserting additional logarithmic factors in the gauge leads to
slower divergent weights and gives the following extension.

\begin{proposition}\label{iterated}
Let \(L_1(r)=\log(1/r)\) and
\(L_{m+1}(r)=\log L_m(r)\).
For every fixed integer \(k\geq2\),
$$
\mathcal P^{h_k}(S)=0,
\qquad h_k(r)=r\prod_{m=1}^kL_m(r)^2,
$$
\end{proposition}

\begin{proof}
Repeat the proof with
$$
F(n)=n^2\prod_{m=1}^{k-1}(\log_m n)^2.
$$
Here \(\log_1 n=\log n\) and \(\log_{m+1}n=\log(\log_m n)\).
For large \(n\), \(F\) increases and \(F(Mj+c)\asymp F(j)\).
The reciprocals \(1/(j\prod_{m=1}^{k-1}\log_m j)\) decrease
and have divergent sum by the integral test. Repeating the
summation by parts in Lemma \ref{divergence} with these explicit
reciprocals, and applying Lemma \ref{trace}, gives the same contradiction.
\end{proof}

\begin{thebibliography}{9}
\small\raggedright
\bibitem{CKN}
L. Caffarelli, R. Kohn, and L. Nirenberg,
\emph{Partial regularity of suitable weak solutions of the Navier--Stokes
equations}, Comm. Pure Appl. Math. \textbf{35} (1982), 771--831.
\href{https://doi.org/10.1002/cpa.3160350604}{doi:10.1002/cpa.3160350604}.

\bibitem{GP}
C. Guevara and N. C. Phuc,
\emph{Local energy bounds and epsilon-regularity criteria for the 3D
Navier--Stokes system}, Calc. Var. Partial Differential Equations
\textbf{56} (2017), article 68.
\href{https://www.math.lsu.edu/~pcnguyen/papers_html/Guevara-Phuc-12-29-16.pdf}
{Author-hosted manuscript}.

\bibitem{LR}
Z. Lei and X. Ren,
\emph{Quantitative partial regularity of the Navier--Stokes equations and
applications}, Adv. Math. \textbf{445} (2024), 109654.
\href{https://doi.org/10.1016/j.aim.2024.109654}{doi:10.1016/j.aim.2024.109654}.

\bibitem{GKT}
S. Gustafson, K. Kang, and T.-P. Tsai,
\emph{Interior regularity criteria for suitable weak solutions of the
Navier--Stokes equations}, Comm. Math. Phys. \textbf{273} (2007), 161--176.
\href{https://doi.org/10.1007/s00220-007-0214-6}{doi:10.1007/s00220-007-0214-6}.

\bibitem{BP}
C. J. Bishop and Y. Peres,
\emph{Fractals in Probability and Analysis}, Cambridge University Press,
2017. Lemma 3.1.1 (Frostman's lemma for a gauge).
\href{https://www.math.stonybrook.edu/~bishop/classes/math324.F15/book1Dec15.pdf}
{Author-hosted manuscript}.

\bibitem{HK}
T. Hyt\"onen and A. Kairema,
\emph{Systems of dyadic cubes in a doubling metric space},
Colloq. Math. \textbf{126} (2012), 1--33.
\href{https://doi.org/10.4064/cm126-1-1}{doi:10.4064/cm126-1-1}.
\end{thebibliography}
\end{document}
